\documentclass[11pt]{article}
\usepackage[margin=1in]{geometry}
\usepackage{amsmath,amssymb,amsthm}
\usepackage{xurl}
\usepackage{hyperref}
\sloppy
\let\oldus\_ \renewcommand{\_}{\oldus\allowbreak}
\newtheorem{theorem}{Theorem}
\newtheorem{lemma}[theorem]{Lemma}
\theoremstyle{definition}
\newtheorem{definition}[theorem]{Definition}

\title{Conjecture 00000001144 is false:\\ a hyperbolic Weyl group with growth rate below 2}
\author{Jackmeson1}
\date{}

\begin{document}
\maketitle

\section{The conjecture and the reading}

\textbf{Statement (English).} ``Conjecture: The growth rate of hyperbolic-type Weyl groups of
Kac--Moody algebras is 2 (hyperbolic Weyl growth).'' The Chinese text says the same: the growth
rate of the Weyl groups of Kac--Moody algebras of hyperbolic type is 2.

\textbf{Reading refuted.} For every hyperbolic generalized Cartan matrix (GCM) $A$, the Weyl
group $W(A)$, with its standard generating set $S=\{r_1,\dots,r_n\}$ of fundamental
reflections, has exponential growth rate $2$. We refute three versions of ``growth rate'': the
$\limsup$ of $a_k^{1/k}$, where $a_k$ is the number of elements of word length exactly $k$; the
$\limsup$ of $|B_k|^{1/k}$, where $B_k$ is the ball of radius $k$; and the limit of either
sequence. One explicit hyperbolic GCM, the matrix of the Feingold--Frenkel algebra
$\mathcal F=A_1^{++}$, has both $\limsup$s at most $19/10$; hence any limit of either sequence,
if it exists, is at most $19/10$, and in particular neither sequence tends to $2$. (We do not
prove that the limits exist; non-convergence to $2$ already refutes the claimed value.)

\textbf{Readings not refuted.} (i) The existential reading, that \emph{some} hyperbolic Weyl
group has growth rate $2$, is true. The rank-3 GCM with all off-diagonal entries $-2$ is
hyperbolic; numerical evidence (BFS up to a finite length, not a proof) suggests that its
Weyl group has $3\cdot 2^{k-1}$ elements of length $k$ for $k\ge 1$. (ii) For the readings ``the
growth rate is at most $2$'' and ``the supremum is $2$'' we have only numerical evidence against
them: for the simply laced rank-4 GCM whose diagram is the complete graph $K_4$ (hyperbolic),
finite BFS gives sphere ratios near $2.303$. A finite BFS proves neither asymptotic convergence
nor an all-length formula; these remarks are evidence only and are not used in, or proved by,
the Lean disproof.

\section{Definitions (as formalized)}

\begin{definition}[GCM, indecomposable]
A GCM is an integer matrix $A=(a_{ij})$ with $a_{ii}=2$, $a_{ij}\le 0$ for $i\ne j$, and
$a_{ij}=0\iff a_{ji}=0$ (Lean: \texttt{GCM}). The matrix $A$ is indecomposable if no nonempty
proper subset $J$ of the indices has $a_{ij}=0$ for all $i\in J$ and $j\notin J$ (Lean:
\texttt{Indecomposable}). Both definitions follow Wikipedia [W1].
\end{definition}

\begin{definition}[Hyperbolic type, symmetrizable case]
We use [W2]. Let $A=DS$ with $D$ diagonal with positive integer entries and $S$ symmetric.
Then $A$ is hyperbolic if it is indecomposable, $S$ is indefinite, and for each proper subset
$J$ of the indices the principal submatrix $S_J$ is positive definite or positive
semidefinite. Lean (\texttt{IsHyperbolic}) asks for: indecomposability; some $d$ with $d_i>0$
and a real symmetric $S$ with $a_{ij}=d_iS_{ij}$; vectors $x,y$ with $x^TSx<0<y^TSy$; and
$x^TSx\ge 0$ for every $x$ supported on a proper subset $J$. The last condition says exactly
that $S_J$ is positive semidefinite.
\end{definition}

The text of [C] gives Kac's general definition: $A$ is hyperbolic if it is neither of finite
nor of affine type, but every proper connected subdiagram is of finite or affine type. [W1]
gives the minor criteria: finite type if all principal minors are positive, affine type if the
proper principal minors are positive and $\det A=0$. For our matrix the Lean lemma
\texttt{FF\_minors} records the data for this definition as well. It gives $\det A=-2$, so $A$ is
neither finite nor affine. The connected proper subdiagrams are the three singletons ($A_1$),
$\{0,1\}$ (minors $2,3$: finite, $A_2$) and $\{1,2\}$ (minors $2$ and $\det=0$: affine,
$A_1^{(1)}$). The subset $\{0,2\}$ is disconnected because $a_{02}=0$.

\begin{definition}[Weyl group, word length, growth rate]
The reflection $r_i$ acts on the root lattice $\bigoplus_j \mathbb Z\alpha_j$ by
$r_i(\alpha_j)=\alpha_j-a_{ij}\alpha_i$. In Lean, \texttt{reflMat} is the integer matrix of
$r_i$ in the basis of simple roots, and \texttt{refl} is the same map as a unit of
$M_n(\mathbb Z)$; we prove $r_i^2=1$. The Weyl group \texttt{weyl} is the subgroup of
$GL_n(\mathbb Z)$ generated by the $r_i$. Following [FKN], the Weyl group associated with a
hyperbolic Cartan matrix is the group generated by the simple reflections. FKN write the
reflections as $w_I(X)=X-\frac{2(X,\alpha_I)}{(\alpha_I,\alpha_I)}\alpha_I$. Our matrix is
symmetric, so transposition conventions do not matter, and $\det A\neq 0$, so the simple roots
are linearly independent. The ball \texttt{ball C k} is the set of products of at most $k$
generators; since $r_i=r_i^{-1}$, words without inverses suffice. The sphere \texttt{sphere C k}
is the set of elements in \texttt{ball C k} and in no smaller ball, i.e.\ of word length exactly
$k$. Theorem \texttt{weyl\_eq\_iUnion\_ball} proves $W(A)=\bigcup_k B_k$. The growth rate is
\texttt{growthRate} $=\limsup_k |{\rm sphere}_k|^{1/k}$. By Cauchy--Hadamard this is the
inverse of the radius of convergence of the growth series $\sum_k a_k t^k$, which is how
[KP] defines it. We also define \texttt{ballGrowthRate} $=\limsup_k |B_k|^{1/k}$, with $B_k$
the ball of [W3].
\end{definition}

Two conventions: \texttt{Set.ncard} is applied only to sets we prove finite, and the
$\limsup$ in $\mathbb R$ is applied only to sequences we prove eventually bounded above and
bounded below by $0$, so neither takes a junk value.

\section{The counterexample}

The Feingold--Frenkel algebra is $\mathcal F=A_1^{++}$. According to [FKN], the discovery
that its Weyl group is $PGL_2(\mathbb Z)$ was the first vital step in Feingold and Frenkel's
work. In the order $\alpha_{-1},\alpha_0,\alpha_1$ (Lean indices $0,1,2$) its GCM is
\[
A=\begin{pmatrix}2&-1&0\\-1&2&-2\\0&-2&2\end{pmatrix}\qquad(\texttt{ffA}).
\]
This is the over-extension of $A_1^{(1)}$: the new node $\alpha_{-1}$ is joined to the
affine node $\alpha_0$ by a single edge. The Lean proof does not depend on the identification: hyperbolicity is proved directly.

\begin{lemma}[\texttt{FF\_hyperbolic}]
$A$ is hyperbolic. Take $D=I$ and $S=A$. The form is
$Q(x)=2x_0^2+2x_1^2+2x_2^2-2x_0x_1-4x_1x_2$. It satisfies $Q(1,2,2)=-2<0$ and $Q(1,0,0)=2>0$.
If $x_0=0$ then $Q=2(x_1-x_2)^2$; if $x_1=0$ then $Q=2x_0^2+2x_2^2$; if $x_2=0$ then
$Q=x_0^2+x_1^2+(x_0-x_1)^2$. Indecomposability is checked by \texttt{decide}.
\end{lemma}

\begin{lemma}[\texttt{FF\_comm}, \texttt{refl\_mul\_self}]
$r_i^2=1$ for all $i$, and $r_0r_2=r_2r_0$.
\end{lemma}

Call a word reduced if it contains neither $ii$ nor the factor $2\,0$ (Lean: chains for the
relation \texttt{okFF}).

\begin{lemma}[\texttt{ff\_prepend}, \texttt{ff\_reduce}]
Every word $l$ has a reduced word $w$ with $|w|\le|l|$ and the same product in $W(A)$.
\end{lemma}
\begin{proof}
Induct on $l$ and prepend one letter $x$ to a reduced word $w=y\,v$. If $xy$ is allowed, take
$x\,y\,v$. If $x=y$, cancel. Otherwise $x=2$ and $y=0$, and $2\,0\,v=0\,2\,v$. If $v$ starts
with $2$, cancel to $0\,v'$; this is reduced, because the letter after $2$ in $v$ is neither
$0$ nor $2$. If $v$ starts with $1$, take $0\,2\,1\,v'$. If $v$ is empty, take $0\,2$.
\end{proof}

\begin{lemma}[\texttt{ff\_follow\_card}, \texttt{ff\_ball\_card}]
The number of reduced words of length $m+1$ that follow a letter $a$ is at most
$v_a(5/3)^{m+1}$, with $v=(5,5,3)$. The induction uses $3(v_1+v_2)\le 5v_0$,
$3(v_0+v_2)\le 5v_1$ and $3v_1\le 5v_2$. Hence at most $13(5/3)^j$ reduced words have length
$j$, and $|B_k|\le 13(k+1)(5/3)^k$; in particular $B_k$ is finite.
\end{lemma}

\begin{theorem}[\texttt{FF\_rate}, \texttt{disproof}]
For all large $k$, both $|{\rm sphere}_k|^{1/k}$ and $|B_k|^{1/k}$ are at most $19/10$. Indeed
$13(k+1)\le(57/50)^k$ eventually, by \texttt{tendsto\_pow\_const\_div\_const\_pow\_of\_one\_lt},
and $(57/50)(5/3)=19/10$. Hence \texttt{growthRate FF} $\le 19/10$ and
\texttt{ballGrowthRate FF} $\le 19/10$; neither sequence of $k$-th roots tends to $2$; and
\[
\neg\,\bigl(\forall n\ \forall C:\texttt{GCM }n,\ \texttt{IsHyperbolic }C\to
\texttt{growthRate }C=2\bigr).
\]
\end{theorem}

Remarks, not used in Lean. Breadth-first search on the reflection matrices up to a finite length
gives sphere ratios close to $1.3247\ldots$, the real root of $x^3-x-1$ (numerical evidence
only; a finite search does not prove convergence). This is consistent with $W\cong
PGL_2(\mathbb Z)\cong T(2,3,\infty)$, since $(r_0r_1)^3=1$ also holds. The rank-2 hyperbolic
GCM $\begin{pmatrix}2&-3\\-3&2\end{pmatrix}$, of which [W2] says there are infinitely many,
has the infinite dihedral group as Weyl group: its spheres have size $2$, so its growth rate
is $1$.

\section{Lean formalization and axiom audit}

File \texttt{lean/Conjecture1144/Basic.lean} (414 lines, \texttt{import Mathlib} only, Lean
v4.33.1 with Mathlib v4.33.1). Main theorem:
\begin{verbatim}
theorem disproof :
    IsHyperbolic FF /\ FF.A.det = -2 /\
    (weyl FF : Set _) = (Union k, ball FF k) /\
    RateBound FF /\ growthRate FF != 2 /\ ballGrowthRate FF != 2 /\
    ~ (forall (n : Nat) (C : GCM n), IsHyperbolic C -> growthRate C = 2)
\end{verbatim}
(ASCII transliteration of the Lean statement.) \texttt{RateBound FF} consists of: every ball is finite; both $k$-th-root sequences are
eventually at most $19/10$; both $\limsup$s are at most $19/10$; and neither sequence tends to
$2$. The command \texttt{\#print axioms C1144.disproof} reports only \texttt{propext},
\texttt{Classical.choice} and \texttt{Quot.sound}. There is no \texttt{sorry} and no
\texttt{native\_decide}; \texttt{decide} is used only on finite statements about explicit
$3\times 3$ integer data.

\begin{thebibliography}{9}
\bibitem[W1]{W1} Wikipedia, ``Cartan matrix'', section on Lie algebras and classification.
\url{https://en.wikipedia.org/wiki/Cartan_matrix}
\bibitem[W2]{W2} Wikipedia, ``Kac--Moody algebra''.
\url{https://en.wikipedia.org/wiki/Kac%E2%80%93Moody_algebra}
\bibitem[W3]{W3} Wikipedia, ``Growth rate (group theory)''.
\url{https://en.wikipedia.org/wiki/Growth_rate_(group_theory)}
\bibitem[C]{C} L. Carbone, S. Chung, L. Cobbs, R. McRae, D. Nandi, Y. Naqvi, D. Penta,
\emph{Classification of hyperbolic Dynkin diagrams, root lengths and Weyl group orbits},
arXiv:1003.0564. \url{https://arxiv.org/abs/1003.0564}
\bibitem[FKN]{FKN} A. J. Feingold, A. Kleinschmidt, H. Nicolai, \emph{Hyperbolic Weyl groups and
the four normed division algebras}, arXiv:0805.3018. \url{https://arxiv.org/abs/0805.3018}
\bibitem[KP]{KP} R. Kellerhals, G. Perren, \emph{On the growth of cocompact hyperbolic Coxeter
groups}, arXiv:0910.4103. \url{https://arxiv.org/abs/0910.4103}
\end{thebibliography}

\end{document}
